% =============================================================================
\section{Cavidades en un dieléctrico}\label{sec:cavidades}
% =============================================================================

Dentro de un dieléctrico el campo macroscópico $\vE$ es un promedio: no podemos introducir una carga de
prueba entre las moléculas para medirlo. Lo que sí podemos hacer es practicar una pequeña cavidad y medir
el campo en su interior. La pregunta es si ese campo coincide con el del medio, y la respuesta es que
depende de la forma de la cavidad. Dos formas extremas, la aguja y la lenteja, dan respectivamente $\vE$ y
$\vD/\eo$, y con ello una definición operacional de ambos campos~\cite{DFA-teoria,Landau}.

\subsection{Planteamiento: el principio de superposición}\label{sec:cav-superposicion}

Sea un dieléctrico con campos $\vE_d$, $\vD_d$ y $\vP_d$, uniformes a la escala de la cavidad que vamos a practicar, y sea
$V_c$ la región de la cavidad, muy pequeña frente a las distancias en que varían los campos.

\begin{proposicion}[Superposición para cavidades]\label{prop:cav-superposicion}
Si la polarización del medio fuera de la cavidad no se altera al practicarla, el campo en la cavidad es
\begin{equation}\label{eq:cav-superposicion}
  \vEcav=\vE_d+\vE_{-\vP_d},
\end{equation}
donde $\vE_{-\vP_d}$ es el campo que crearía un cuerpo con la forma de la cavidad y polarización uniforme $-\vP_d$.
\end{proposicion}

\begin{proof}
La polarización del medio con cavidad es $\vP=\vP_d-\vP_d\,\mathbb 1_{V_c}$, donde $\mathbb 1_{V_c}$ vale uno en la cavidad y cero fuera.
El potencial \eqref{eq:phi-P} depende linealmente de $\vP$, así que el campo es la suma del creado por el medio macizo,
$\vE_d$, y del creado por la polarización $-\vP_d$ restringida a $V_c$ (\cref{fig:cavidades}\,(a)). Las cargas libres, si las hay,
no cambian.
\end{proof}

\begin{observacion}[Sobre la hipótesis]\label{obs:cav-hipotesis}
En un electrete la hipótesis es exacta: su polarización no depende del campo. En un medio \LHI{} la cavidad modifica
ligeramente el campo y la polarización en sus proximidades, y la \cref{prop:cav-superposicion} es solo aproximada. Los
resultados de las cavidades acicular y lenticular, sin embargo, son exactos en el límite de forma extrema, porque se
deducen directamente de las condiciones de frontera, como veremos.
\end{observacion}

% Figura: cavidades en un dieléctrico. Principio de superposición, cavidad
% acicular y cavidad lenticular, con los perfiles axiales que las justifican.
\begin{figure}[t]
\centering
\begin{subfigure}[b]{\textwidth}\centering
\begin{tikzpicture}[scale=0.92]
  \tikzset{bloque/.style={fill=dielec, draw=texto!45!fondo, line width=0.4pt, rounded corners=2pt}}
  \tikzset{aguja/.style={fill=fondo, draw=texto, line width=0.5pt, rounded corners=1.4pt}}
  % dieléctrico con cavidad
  \draw[bloque] (0,0) rectangle (3.0,1.9);
  \draw[aguja] (0.95,0.7) rectangle (2.05,0.86);
  \draw[vecP] (1.05,1.38) -- (1.95,1.38) node[above, midway, etiq] {$\vP_d$};
  \node[etiq] at (3.55,0.95) {$=$};
  % dieléctrico macizo
  \draw[bloque] (4.1,0) rectangle (7.1,1.9);
  \draw[vecP] (5.15,1.38) -- (6.05,1.38) node[above, midway, etiq] {$\vP_d$};
  \node[etiq] at (7.65,0.95) {$+$};
  % «anticavidad»: cilindro con -P_d
  \draw[aguja, fill=dielecB] (8.25,0.7) rectangle (9.75,0.86);
  \draw[vecP] (9.55,0.78) -- (8.45,0.78);
  \node[etiq, anchor=north] at (9.0,0.6) {$-\vP_d$};
\end{tikzpicture}
\caption{Superposición: dieléctrico con cavidad $=$ dieléctrico macizo $+$ cuerpo con la forma de la cavidad y polarización $-\vP_d$}
\end{subfigure}\\[1.0em]
\begin{subfigure}[b]{0.235\textwidth}\centering
\begin{tikzpicture}[scale=0.86]
  \fill[dielec, rounded corners=2pt] (0,0) rectangle (3.0,2.4);
  \draw[vecE] (0.25,2.0) -- (1.15,2.0) node[above, midway, etiq] {$\vE_d$};
  \draw[vecP] (1.8,2.0) -- (2.7,2.0) node[above, midway, etiq] {$\vP_d$};
  \draw[fill=fondo, draw=texto, line width=0.5pt, rounded corners=1.4pt] (0.55,0.62) rectangle (2.45,0.8);
  \draw[vecn] (1.5,0.8) -- (1.5,1.42) node[right, etiq] {$\nn$};
  \draw[vecn] (1.5,0.8) -- (2.1,0.8);
  \node[etiq, anchor=south] at (2.05,0.8) {$\tu$};
  \node[etiq, anchor=north] at (1.5,0.55) {$\vEcav=\vE_d$};
\end{tikzpicture}
\caption{Acicular}
\end{subfigure}\hfill
\begin{subfigure}[b]{0.24\textwidth}\centering
\begin{tikzpicture}
\begin{axis}[perfil, width=4.25cm, height=3.6cm, ymin=-0.75, ymax=1.15, ytick={-0.5,0.5,1}]
  \curvaPeje \curvaEeje{0.1}
  \node[etiqf, font=\tiny, anchor=north west] at (rel axis cs:0.02,0.98) {$a/\ell=0{,}1$};
\end{axis}
\end{tikzpicture}
\caption{Cilindro largo}
\end{subfigure}\hfill
\begin{subfigure}[b]{0.235\textwidth}\centering
\begin{tikzpicture}[scale=0.86]
  \fill[dielec, rounded corners=2pt] (0,0) rectangle (3.0,2.4);
  \draw[vecE] (0.25,2.0) -- (1.15,2.0) node[above, midway, etiq] {$\vE_d$};
  \draw[vecP] (1.8,2.0) -- (2.7,2.0) node[above, midway, etiq] {$\vP_d$};
  \draw[fill=fondo, draw=texto, line width=0.5pt] (1.1,0.95) ellipse (0.09 and 0.75);
  \draw[vecn] (1.19,0.95) -- (1.8,0.95) node[above, etiq] {$\nn$};
  \draw[vecn] (1.19,0.95) -- (1.19,1.5) node[right, etiq] {$\tu$};
  \node[etiq, anchor=west] at (1.3,0.42) {$\vDcav=\vD_d$};
\end{tikzpicture}
\caption{Lenticular}
\end{subfigure}\hfill
\begin{subfigure}[b]{0.24\textwidth}\centering
\begin{tikzpicture}
\begin{axis}[perfil, width=4.25cm, height=3.6cm, ymin=-1.15, ymax=1.15, ytick={-1,1}]
  \curvaPeje \curvaEeje{50}
  \node[etiqf, font=\tiny, anchor=north west] at (rel axis cs:0.02,0.98) {$a/\ell=50$};
\end{axis}
\end{tikzpicture}
\caption{Disco delgado}
\end{subfigure}
\caption{Cavidades en un dieléctrico polarizado ($\vE_d$, $\vP_d$, $\vD_d$ son los campos del medio lejos de la
cavidad). (a) Por superposición, basta sumar al campo del medio el de un cuerpo con la forma de la
cavidad y polarización $-\vP_d$. (b) En una cavidad acicular paralela al campo domina la superficie
lateral y se conserva la componente tangencial: $\vEcav=\vE_d$. (c) Perfil axial de $E$ (línea continua,
unidades $P/\eo$) y de $P$ (trazos) de un cilindro polarizado largo: dentro, lejos de los extremos, $E\simeq0$.
(d) En una cavidad lenticular perpendicular al campo dominan las caras y se conserva la componente
normal de $\vD$: $\vDcav=\vD_d$. (e) En un disco delgado, $E\simeq-P/\eo$ dentro; con $-\vP_d$ esto aporta $+\vP_d/\eo$.}
\label{fig:cavidades}
\end{figure}


\subsection{Cavidad acicular}\label{sec:acicular}

\begin{definicion}[Cavidad acicular]\label{def:acicular}
Una \emph{cavidad acicular} es un cilindro de radio $a$ y longitud $2\ell$ con $a\ll\ell$ (una aguja), cuyo eje es paralelo a $\vE_d$.
\end{definicion}

\begin{proposicion}[Campo en una cavidad acicular]\label{prop:acicular}
En el centro de una cavidad acicular, en el límite $a/\ell\to0$,
\begin{equation}\label{eq:acicular}
  \boxed{\;\vEcav=\vE_d\;},\qquad \vDcav=\eo\vE_d=\vD_d-\vP_d .
\end{equation}
\end{proposicion}

\begin{proof}
\emph{Por condiciones de frontera.} La superficie de la aguja es casi toda lateral, y en ella $\vE_d$ es tangente. La componente
tangencial de $\vE$ es continua (\cref{teo:frontera}), luego justo dentro de la pared lateral $\vEcav=\vE_d$. Las caras de los
extremos tienen área $\pi a^2$, despreciable, y su influencia en el centro desaparece en el límite.

\emph{Por superposición.} El cuerpo de la \cref{prop:cav-superposicion} es un cilindro con polarización axial $-\vP_d$, que equivale a dos
discos de radio $a$ con densidades $\pm P_d$ en $z=\mp\ell$. En el centro, por el resultado del \cref{prob:cilindro-polarizado} con
$P\to-P_d$ y $z=0$,
\[
  \vE_{-\vP_d}(0)=\frac{\vP_d}{\eo}\left(1-\frac{\ell}{\sqrt{a^2+\ell^2}}\right)=\frac{\vP_d}{\eo}\left(\frac{a^2}{2\ell^2}+\mathcal O\big(a^4/\ell^4\big)\right)
  \xrightarrow{\;a/\ell\to0\;}\vcero ,
\]
como muestra el perfil del cilindro largo, en el que $E\simeq0$ lejos de los extremos (\cref{fig:cavidades}\,(c)).
\end{proof}

\subsection{Cavidad lenticular}\label{sec:lenticular}

\begin{definicion}[Cavidad lenticular]\label{def:lenticular}
Una \emph{cavidad lenticular} es un cilindro de radio $a$ y espesor $2\ell$ con $\ell\ll a$ (una lenteja o disco), cuyas caras son
perpendiculares a $\vE_d$.
\end{definicion}

\begin{proposicion}[Campo en una cavidad lenticular]\label{prop:lenticular}
En el centro de una cavidad lenticular, en el límite $\ell/a\to0$,
\begin{equation}\label{eq:lenticular}
  \boxed{\;\vDcav=\vD_d\;},\qquad \vEcav=\frac{\vD_d}{\eo}=\vE_d+\frac{\vP_d}{\eo} .
\end{equation}
\end{proposicion}

\begin{proof}
\emph{Por condiciones de frontera.} Ahora la superficie es casi toda cara, perpendicular a $\vE_d$ y por tanto a $\vD_d$. No hay carga libre
en ella, y la componente normal de $\vD$ es continua (\cref{teo:frontera}): justo dentro de la cara, $\vDcav=\vD_d$. El borde lateral, de
área $4\pi a\ell$, es despreciable frente a las caras en el límite.

\emph{Por superposición.} El cuerpo es ahora un disco delgado con $-\vP_d$. En su centro, con la misma fórmula y $\ell/a\to0$,
\[
  \vE_{-\vP_d}(0)=\frac{\vP_d}{\eo}\left(1-\frac{\ell}{\sqrt{a^2+\ell^2}}\right)\xrightarrow{\;\ell/a\to0\;}\frac{\vP_d}{\eo},
\]
el campo de un condensador plano con densidades $\pm P_d$, como en el perfil del disco delgado (\cref{fig:cavidades}\,(e)). Sumando
$\vE_d$ se obtiene \eqref{eq:lenticular}.
\end{proof}

\subsection{Interpretación: definición operacional de $\vE$ y $\vD$}\label{sec:cav-interpretacion}

Las dos proposiciones anteriores, debidas esencialmente a Kelvin, permiten definir los campos macroscópicos mediante
medidas: \emph{$\vE$ es la fuerza por unidad de carga que actúa sobre una carga de prueba en una cavidad acicular
paralela al campo, y $\vD/\eo$ es la que actúa en una cavidad lenticular perpendicular a él}. La diferencia entre ambas
es $\vP/\eo$: es el campo de las cargas ligadas de las caras del disco, que la aguja evita y la lenteja no.

\begin{observacion}[La cavidad esférica y el campo de Lorentz]\label{obs:lorentz}
Para una cavidad esférica, el cuerpo de la \cref{prop:cav-superposicion} es una esfera con $-\vP_d$, que crea en su interior el campo
uniforme $+\vP_d/3\eo$ (\cref{prob:esfera-polarizada}). Por tanto,
\[
  \vEcav=\vE_d+\frac{\vP_d}{3\eo},
\]
a medio camino entre la aguja y la lenteja. Es el llamado \emph{campo local de Lorentz}: una estimación del campo que actúa
sobre una molécula, que «ve» a sus vecinas como un medio continuo a partir de cierta distancia. Si cada molécula
responde con $\vp=\alpha\vE_{\mathrm{loc}}$ y hay $N$ por unidad de volumen, $\vP=N\alpha\,(\vE+\vP/3\eo)$, y despejando $\chi=P/(\eo E)$ se obtiene
la relación de Clausius--Mossotti~\cite{Griffiths,Jackson},
\[
  \frac{\er-1}{\er+2}=\frac{N\alpha}{3\eo},
\]
que liga la polarizabilidad microscópica de la \cref{sec:microscopico} con la permitividad macroscópica de la
\cref{sec:lineales}.
\end{observacion}

\begin{observacion}[Cavidad elipsoidal]\label{obs:elipsoide}
Más en general, una cavidad con forma de elipsoide de revolución cuyo eje es paralelo a $\vE_d$ tiene en su interior un campo
uniforme $\vEcav=\vE_d+N\,\vP_d/\eo$, donde el \emph{factor de despolarización} $N\in[0,1]$ solo depende de la relación de aspecto~\cite{Osborn,Landau}.
Vale $N\to0$ para la aguja, $N=1/3$ para la esfera y $N\to1$ para la lenteja: las tres cavidades de esta sección son puntos de
una misma curva (\cref{fig:despolarizacion}).
\end{observacion}

% Figura: factor de despolarización N de un esferoide de revolución (eje paralelo
% al campo) en función de su relación de aspecto m = c/a.  Datos: scripts/lineas_campo.py
\begin{figure}[t]
\centering
\begin{tikzpicture}
\begin{axis}[width=10.5cm, height=5.4cm, xmode=log, axis lines=left, xmin=0.01, xmax=100,
  ymin=0, ymax=1.08, grid=major, xlabel={relación de aspecto $m=c/a$ (eje paralelo / eje transversal)},
  ylabel={$N(m)$}, ytick={0,0.3333,1}, yticklabels={$0$,$\tfrac13$,$1$},
  xtick={0.01,0.1,1,10,100}, xticklabels={$10^{-2}$,$10^{-1}$,$1$,$10$,$10^{2}$}, clip=false]
  \addplot[curvaE, line width=1.1pt] table[x=m, y=N] {datos/despolarizacion.dat};
  \addplot[only marks, mark=*, mark size=1.8pt, color=acentoC] coordinates {(1,0.3333)};
  % iconos de forma
  \draw[texto, fill=dielecB, line width=0.4pt] (axis cs:0.016,0.42) ellipse [x radius=0.05cm, y radius=0.36cm];
  \node[etiqt, anchor=west] at (axis cs:0.022,0.42) {lenteja: $N\to1$,\; $\vDcav=\vD_d$};
  \draw[texto, fill=dielecB, line width=0.4pt] (axis cs:1,0.58) circle (0.17cm);
  \node[etiqt, anchor=west] at (axis cs:1.35,0.58) {esfera: $N=\tfrac13$ (campo de Lorentz)};
  \draw[texto, fill=dielecB, line width=0.4pt] (axis cs:30,0.25) ellipse [x radius=0.42cm, y radius=0.05cm];
  \node[etiqt, anchor=south] at (axis cs:30,0.29) {aguja: $N\to0$, $\vEcav=\vE_d$};
  \draw[vecE] (axis cs:0.013,0.15) -- (axis cs:0.04,0.15);
  \node[etiqt, anchor=west] at (axis cs:0.042,0.15) {$\vE_d$};
\end{axis}
\end{tikzpicture}
\caption{Factor de despolarización $N$ de una cavidad con forma de esferoide de revolución cuyo eje,
de semieje $c$, es paralelo a $\vE_d$ (el otro semieje es $a$). El campo en la cavidad es
$\vEcav=\vE_d+N\,\vP_d/\eo$, que interpola entre los dos casos límite del texto: la aguja ($m\to\infty$) y la
lenteja ($m\to0$). La esfera da el campo local de Lorentz, $\vE_d+\vP_d/3\eo$.}
\label{fig:despolarizacion}
\end{figure}

